% problem-000181 3 金属晶体与氧化物结构
\section*{3. 金属晶体与氧化物结构}
随着仪器和表征⽅法的发展，各种体系特别是复杂体系热⼒学参数的测定⽅法越来越多样化。例如可以利⽤核磁共振(NMR)波谱确定复杂体系的酸解离常数，下⾯给出⼀个例⼦：

（图：）

此为(ATP)H3−与(ATP)4−之间的质⼦解离-结合平衡过程。为书写⽅便，(ATP)H3−简写为 $\mathrm { H A ^ { 3 - } }$ $\left( \mathrm { A T P } \right) ^ { 4 - }$ 简写为 $\mathrm { A } ^ { 4 - }$ ，酸解离常数为 $K _ { a } ,$ 。简化表达如下：

$
\mathrm{HA} ^ {3 -} \rightarrow \mathrm{H} ^ {+} + \mathrm{A} ^ {4 -}
$

随着pH变化，磷的化学环境发⽣变化，NMR中的化学位移随之变化。下图给出γ位31P的化学位移测定值随pH的变化关系。由于溶液中质⼦交换速度很快，实际测得的化学位移 $( \delta _ { \mathrm { o b s } } )$ 是质⼦化形式化学位移 $( \delta _ { \mathrm { H A } ^ { 3 - } } )$ 和脱质⼦形式化学位移 $\left( { \delta _ { { \tt A } ^ { 4 - } } } \right)$ 的加权平均值：

$
\delta_ {\mathrm{obs}} = \delta_ {\mathrm {HA^ {3 - }}} \chi_ {\mathrm {HA^ {3 -}}} + \delta_ {\mathrm {A^ {4 - }}} \chi_ {\mathrm {A^ {4 - }}}
$

式中， $\chi _ { \mathrm { H A } ^ { 3 - } } \Re \mathrm { H } \chi _ { \mathrm { A } ^ { 4 - } }$ 分别表⽰H $\mathrm { A ^ { 3 - } }$ −和 $\mathrm { A } ^ { 4 ^ { - } }$ 的⽐例。

（图：）
pH

\subsection*{3-1}
写出质⼦化形式化学位移 $( \delta _ { \mathrm { H A } ^ { 3 - } }$ _{"})和脱质⼦形式化学位移 $\left( { \delta _ { \mathrm { { A } } ^ { 4 - } } } \right)$ 的值（要求读到⼩数点后⼀位）。

\subsection*{3-2}
推导出 $\mathsf { p } K _ { a }$ 的表达式，只可包含pH和所有相关的化学位移参数。

\subsection*{3-3}
从上图中合理取值，计算酸解离常数 $K _ { a }$ 。
